\newpage
\section*{Solucionario completo}

\noindent Cada respuesta incluye la clave oficial y su justificación analítica o matemática detallada.

\subsection*{Bloque A. Fundamentos teóricos y conceptuales}

\sol{1}{C}{\textbf{Tema:} Valor del dinero en el tiempo. \textbf{Explicación:} El dinero tiene valor temporal por tres factores intrínsecos: la impaciencia racional o preferencia por consumir hoy, la pérdida de poder de compra por inflación y el costo de oportunidad del capital (posibilidad de generar intereses libres de riesgo).}

\sol{2}{A}{\textbf{Tema:} Interés compuesto vs simple. \textbf{Explicación:} En el interés simple el rendimiento se calcula siempre sobre el principal inicial ($VF = VP(1 + i \cdot n)$). En el interés compuesto los intereses devengados se reinvierten sistemáticamente, generando un crecimiento exponencial ($VF = VP(1 + i)^n$).}

\sol{3}{C}{\textbf{Tema:} Tasa Efectiva Anual (TEA). \textbf{Explicación:} La TNA es una tasa de referencia meramente proporcional. La TEA refleja el rendimiento o costo efectivo generado al cabo de un año cuando los intereses se capitalizan $m$ veces al año (mensual, diario, etc.), cumpliéndose siempre $TEA > TNA$ para capitalización subanual.}

\sol{4}{B}{\textbf{Tema:} Ecuación de Fisher exacta. \textbf{Explicación:} La versión exacta de Irving Fisher descuenta el efecto inflacionario del principal y de los intereses: $(1 + i) = (1 + r)(1 + \pi) = 1 + r + \pi + r\pi$. La aproximación lineal $i \approx r + \pi$ ignora el término cruzado $r\pi$, el cual es relevante cuando la inflación es moderada o alta.}

\sol{5}{D}{\textbf{Tema:} Valor presente de una perpetuidad simple. \textbf{Explicación:} Sumando los términos de una serie geométrica infinita decreciente con razón $\frac{1}{1+r}$: $VP = \sum_{t=1}^{\infty} \frac{C}{(1+r)^t} = \frac{C}{1+r} \frac{1}{1 - 1/(1+r)} = \frac{C}{r}$.}

\sol{6}{C}{\textbf{Tema:} Modelo de Gordon-Shapiro. \textbf{Explicación:} Cuando los dividendos o flujos crecen geométricamente a tasa $g$, el valor presente de la serie infinita converge formalmente a $VP = \frac{C_1}{r - g}$, siempre y cuando la tasa de descuento supere a la tasa de crecimiento ($r > g$).}

\sol{7}{B}{\textbf{Tema:} Anualidades anticipadas vs vencidas. \textbf{Explicación:} Al recibirse las cuotas al inicio de cada periodo, cada flujo está descontado por un factor $(1+r)$ menos que en la anualidad ordinaria. Por ende, $VP_{ant} = VP_{venc} \times (1 + r)$.}

\sol{8}{C}{\textbf{Tema:} Pendiente de la curva de rendimiento. \textbf{Explicación:} La teoría de la prima de liquidez (Hicks) explica que ante el mayor riesgo de duración y volatilidad del precio de los activos a largo plazo, los agentes demandan un mayor retorno esperado para inmovilizar su liquidez en horizontes dilatados.}

\sol{9}{A}{\textbf{Tema:} Ecuación fundamental de la contabilidad. \textbf{Explicación:} El Activo Total (totalidad de bienes y derechos controlados por la entidad) está íntegramente financiado por dos fuentes: deuda con terceros ajenos a la empresa (Pasivo) y capital aportado o retenido por los propietarios (Patrimonio Neto).}

\sol{10}{A}{\textbf{Tema:} Prueba ácida vs Razón corriente. \textbf{Explicación:} $Prueba\ \acute{A}cida = \frac{Activo\ Corriente - Inventarios}{Pasivo\ Corriente}$. Permite evaluar si la empresa puede honrar sus deudas exigibles de corto plazo sin depender de la venta o liquidación forzosa de sus existencias.}

\sol{11}{B}{\textbf{Tema:} Capital de Trabajo Neto. \textbf{Explicación:} El CTN mide el margen de seguridad financiera de corto plazo: la porción de los activos circulantes financiada con pasivos de largo plazo o recursos patrimoniales permanentes.}

\sol{12}{A}{\textbf{Tema:} Ratio de apalancamiento D/E. \textbf{Explicación:} El cociente $D/E = Pasivo\ Total / Patrimonio\ Neto$ cuantifica el riesgo financiero: un elevado apalancamiento amplifica los retornos del accionista en periodos de bonanza pero eleva drásticamente la probabilidad de quiebra ante caídas en los ingresos operativos.}

\sol{13}{B}{\textbf{Tema:} Ratio de cobertura de intereses. \textbf{Explicación:} Indica cuántas veces el beneficio generado por las operaciones del negocio cubre la carga fija de intereses de la deuda. Ratios inferiores a 1.5 a 2.0 veces alertan a los bancos sobre alto riesgo de default crediticio.}

\sol{14}{D}{\textbf{Tema:} ROA vs ROE. \textbf{Explicación:} $ROA = \frac{Utilidad\ Neta}{Activo\ Total}$ mide la productividad integral de las inversiones; $ROE = \frac{Utilidad\ Neta}{Patrimonio}$ mide el rendimiento financiero del propietario, el cual se apalanca con el uso de endeudamiento si el costo de la deuda es menor que el retorno de los activos.}

\sol{15}{B}{\textbf{Tema:} Análisis DuPont del ROE. \textbf{Explicación:} $ROE = \frac{Utilidad\ Neta}{Ventas} \times \frac{Ventas}{Activos} \times \frac{Activos}{Patrimonio} = Margen \times Rotaci\acute{o}n \times Apalancamiento$. Permite identificar si la rentabilidad proviene de margen de precios, volumen de ventas o endeudamiento.}

\sol{16}{A}{\textbf{Tema:} Ciclo de conversión de efectivo. \textbf{Explicación:} $CCE = D\acute{\imath}as\ Inventario + D\acute{\imath}as\ Cobranza - D\acute{\imath}as\ Pago$. Un menor CCE reduce los requerimientos de financiamiento de capital de trabajo y eleva la liquidez del negocio.}

\sol{17}{C}{\textbf{Tema:} Flujo de caja libre del proyecto. \textbf{Explicación:} El FCLP evalúa la rentabilidad económica del activo puro. Los intereses y amortizaciones no se restan del flujo porque el costo de financiamiento de la deuda ya está incorporado en la tasa de descuento utilizada para descontarlo ($WACC$).}

\sol{18}{A}{\textbf{Tema:} Tratamiento de costos hundidos. \textbf{Explicación:} La regla de decisión financiera exige evaluar únicamente flujos de caja incrementales futuros ($Marginales$). Los costos pasados e irreversibles son irrelevantes para la decisión prospectiva de seguir o abandonar el proyecto.}

\sol{19}{D}{\textbf{Tema:} Regla de decisión del VAN. \textbf{Explicación:} El $VAN = \sum_{t=1}^n \frac{FC_t}{(1+r)^t} - I_0$ cuantifica el incremento absoluto en la riqueza de la empresa por encima del costo de oportunidad del capital. Todo proyecto con $VAN > 0$ genera valor y debe ejecutarse.}

\sol{20}{B}{\textbf{Tema:} Aditividad de valor del VAN. \textbf{Explicación:} La aditividad de valor permite evaluar proyectos de manera descentralizada o sumar el valor de divisiones de negocio sin distorsiones matemáticas, propiedad que no cumple la TIR ni el Payback.}

\sol{21}{A}{\textbf{Tema:} Definición matemática de la TIR. \textbf{Explicación:} La TIR es la solución $r$ a la ecuación polinómica $\sum_{t=0}^n \frac{FC_t}{(1+r)^t} = 0$. Representa la rentabilidad porcentual interna intrínseca que rinde el capital invertido en el proyecto a lo largo de su vida útil.}

\sol{22}{B}{\textbf{Tema:} Supuesto de reinversión VAN vs TIR. \textbf{Explicación:} Si un proyecto tiene una TIR del 80\%, la fórmula asume que cada flujo recibido puede reinvertirse de inmediato al 80\%, inflando artificialmente el atractivo del proyecto. El VAN es más realista al asumir reinversión al costo de oportunidad de mercado ($WACC$).}

\sol{23}{D}{\textbf{Tema:} TIR múltiple por cambios de signo. \textbf{Explicación:} Cada cambio de signo en el polinomio de flujos de caja genera una raíz real positiva potencial. Al haber múltiples TIR, no existe una tasa única con la cual comparar el costo de capital, debiendo descartarse la TIR a favor del VAN.}

\sol{24}{D}{\textbf{Tema:} Conflicto VAN vs TIR en proyectos excluyentes. \textbf{Explicación:} El objetivo de la corporación es maximizar la riqueza total de los accionistas (valor absoluto en soles), no la tasa porcentual por sol invertido. Ante conflicto entre proyectos excluyentes, la regla del VAN superior manda siempre sobre la TIR.}

\sol{25}{B}{\textbf{Tema:} Ventajas de la TIRM. \textbf{Explicación:} La TIRM combina todos los flujos de entrada en un único Valor Futuro acumulado a la tasa de costo de capital y los desembolsos en un único Valor Presente. Al haber un solo cambio de signo temporal, garantiza una tasa única y sin el sesgo de reinversión a tasas extremas.}

\sol{26}{B}{\textbf{Tema:} Limitaciones del Payback simple. \textbf{Explicación:} El Payback suma flujos nominales sin descontar (tratando un sol del año 5 igual que un sol de hoy) y castiga arbitrariamente proyectos estratégicos de larga vida útil que generan grandes beneficios tras el año de corte establecido.}

\sol{27}{A}{\textbf{Tema:} Costo Anual Equivalente (CAE). \textbf{Explicación:} Cuando dos plantas o bombas de agua tienen distinta duración (ej. 5 años vs 10 años) y deben reponerse indefinidamente, no se pueden comparar sus VAN directos. El CAE anualiza el costo presente total mediante una anualidad uniforme ($CAE = VP(Costos) / A_n^r$), eligiéndose el menor CAE.}

\sol{28}{B}{\textbf{Tema:} Fórmula del WACC corporativo. \textbf{Explicación:} El WACC pondera el costo de cada fuente de financiamiento por su participación en el valor total de la firma ($V$). La deuda incorpora el factor $(1 - T)$ porque el pago de intereses es deducible del impuesto a la renta (escudo fiscal).}

\sol{29}{C}{\textbf{Tema:} Concepto de Beta en el CAPM. \textbf{Explicación:} El inversionista en portafolios diversificados elimina el riesgo no sistemático (idiosincrático). Por tanto, el mercado solo premia el riesgo que no se puede diversificar (riesgo sistemático medido por $\beta = \frac{Cov(R_i, R_m)}{Var(R_m)}$).}

\sol{30}{C}{\textbf{Tema:} Modigliani-Miller Proposición I (sin impuestos). \textbf{Explicación:} Bajo ausencia de impuestos y fricciones, el apalancamiento no genera valor: el mayor riesgo financiero de la deuda eleva el costo de los recursos propios ($K_e$) en la proporción exacta para mantener constante el $WACC$, dejando el valor de la firma inalterado.}

\subsection*{Bloque B. Aplicaciones cuantitativas y financieras}

\sol{31}{D}{\textbf{Tema:} Conversión de TNA mensual a TEA. \textbf{Explicación:} La tasa efectiva mensual es $i_m = 12\% / 12 = 1\% = 0.01$. Capitalizando a 12 meses: $TEA = (1 + 0.01)^{12} - 1 = (1.01)^{12} - 1 \approx 1.126825 - 1 = 12.68\%$.}

\sol{32}{D}{\textbf{Tema:} Cálculo de tasa real de Fisher. \textbf{Explicación:} Aplicando la fórmula exacta: $1 + r = \frac{1 + i}{1 + \pi} = \frac{1.0815}{1.0300} = 1.0500 \implies r = 5.00\%$ anual.}

\sol{33}{A}{\textbf{Tema:} Valor presente de perpetuidad simple. \textbf{Explicación:} $VP = \frac{C}{r} = \frac{50000}{0.08} = 625000$ soles.}

\sol{34}{B}{\textbf{Tema:} Cálculo con modelo de Gordon-Shapiro. \textbf{Explicación:} $VP = \frac{C_1}{r - g} = \frac{30000}{0.09 - 0.03} = \frac{30000}{0.06} = 500000$ soles.}

\sol{35}{B}{\textbf{Tema:} Cálculo de cuota fija de anualidad. \textbf{Explicación:} $C = P \times \left[ \frac{r(1+r)^n}{(1+r)^n - 1} \right] = 100000 \times 0.2637975 = 26379.75$ soles anuales.}

\sol{36}{C}{\textbf{Tema:} Cálculo de razón corriente y prueba ácida. \textbf{Explicación:} $Razón\ Corriente = \frac{160000}{80000} = 2.0$ veces. $Prueba\ \acute{A}cida = \frac{160000 - 60000}{80000} = \frac{100000}{80000} = 1.25$ veces.}

\sol{37}{A}{\textbf{Tema:} Cálculo de D/E ratio. \textbf{Explicación:} Patrimonio $= Activo - Pasivo = 500000 - 300000 = 200000$ soles. $D/E = \frac{Deuda}{Patrimonio} = \frac{300000}{200000} = 1.50$ (o $150\%$).}

\sol{38}{A}{\textbf{Tema:} Cálculo de cobertura de intereses. \textbf{Explicación:} $Cobertura = \frac{EBIT}{Intereses} = \frac{120000}{30000} = 4.00$ veces. La empresa genera 4 soles de beneficio operativo por cada sol que debe pagar en intereses.}

\sol{39}{C}{\textbf{Tema:} Cálculo del ROE con DuPont. \textbf{Explicación:} $ROE = Margen \times Rotaci\acute{o}n \times Apalancamiento = 0.08 \times 1.5 \times 1.50 = 0.18 = 18.0\%$.}

\sol{40}{B}{\textbf{Tema:} Cálculo del ciclo de conversión del efectivo. \textbf{Explicación:} $CCE = D\acute{\imath}as\ Cobro + D\acute{\imath}as\ Inventario - D\acute{\imath}as\ Pago = 45 + 60 - 40 = 65$ días.}

\sol{41}{C}{\textbf{Tema:} Cálculo de Flujo de Caja Libre (FCLP). \textbf{Explicación:} $FCLP = EBIT(1 - T) + Depreciaci\acute{o}n - CAPEX - \Delta CTN = 200000(0.70) + 40000 - 50000 - 10000 = 140000 + 40000 - 60000 = 120000$ soles.}

\sol{42}{B}{\textbf{Tema:} Cálculo del VAN de un proyecto. \textbf{Explicación:} $VAN = -100000 + \frac{40000}{1.10} + \frac{40000}{(1.10)^2} + \frac{40000}{(1.10)^3} = -100000 + 36363.6 + 33057.9 + 30052.6 = -525.9$ soles. Al ser $VAN < 0$, el proyecto debe rechazarse.}

\sol{43}{B}{\textbf{Tema:} Cálculo de TIR de 2 periodos. \textbf{Explicación:} Sea $x = 1/(1+TIR)$. La ecuación es $10000 = 6000x + 6000x^2 \implies 6x^2 + 6x - 10 = 0 \implies 3x^2 + 3x - 5 = 0$. $x = \frac{-3 + \sqrt{9 - 4(3)(-5)}}{6} = \frac{-3 + \sqrt{69}}{6} \approx \frac{-3 + 8.3066}{6} = 0.88443$. Luego $1+TIR = 1/0.88443 \approx 1.13066 \implies TIR \approx 13.07\%$.}

\sol{44}{C}{\textbf{Tema:} Cálculo de la Tasa de Fisher. \textbf{Explicación:} La diferencia de flujos (A - B) es: $t_0 = 0$, $t_1 = 700 - 200 = 500$, $t_2 = 700 - 1300 = -600$. Igualando a cero: $\frac{500}{1+r} - \frac{600}{(1+r)^2} = 0 \implies 500 = \frac{600}{1+r} \implies 1 + r = \frac{600}{500} = 1.20 \implies r_{Fisher} = 20.0\%$.}

\sol{45}{C}{\textbf{Tema:} Cálculo de Payback simple. \textbf{Explicación:} $Payback = \frac{Inversi\acute{o}n}{Flujo\ Anual} = \frac{70000}{25000} = 2.80$ años.}

\sol{46}{D}{\textbf{Tema:} Relación Beneficio/Costo (B/C). \textbf{Explicación:} $B/C = \frac{VP(Beneficios)}{VP(Costos)} = \frac{280000}{200000} = 1.40$. Al ser $B/C > 1$, el proyecto es rentable y genera un $VAN = 280000 - 200000 = +80000$ soles.}

\sol{47}{B}{\textbf{Tema:} Cálculo de Costo Anual Equivalente (CAE). \textbf{Explicación:} $CAE_{Alfa} = \frac{50000}{2.48685} \approx 20105.75$ soles. $CAE_{Beta} = \frac{75000}{3.79079} \approx 19784.80$ soles. Aunque Beta requiere mayor costo inicial total, su mayor durabilidad genera un menor costo anualizado ($19784.8 < 20105.7$).}

\sol{48}{A}{\textbf{Tema:} Cálculo de WACC con impuestos. \textbf{Explicación:} $WACC = K_e (E/V) + K_d (1 - T)(D/V) = 0.12(0.60) + 0.07(1 - 0.30)(0.40) = 0.072 + 0.07(0.70)(0.40) = 0.072 + 0.0196 = 0.0916 = 9.16\%$.}

\sol{49}{C}{\textbf{Tema:} Cálculo de costo de capital con CAPM emergente. \textbf{Explicación:} $K_e = R_f + \beta(R_m - R_f) + RP = 0.04 + 0.75(0.06) + 0.015 = 0.04 + 0.045 + 0.015 = 0.1000 = 10.00\%$.}

\sol{50}{B}{\textbf{Tema:} Cálculo del escudo fiscal de la deuda. \textbf{Explicación:} El ahorro tributario anual en intereses es $T \times (K_d \times D) = 0.30 \times 0.06 \times 3000000 = 54000$ soles. Descontando esta perpetuidad al costo de la deuda $K_d$: $VP = \frac{T \times K_d \times D}{K_d} = T \times D = 0.30 \times 3000000 = 900000$ soles.}

\sol{51}{A}{\textbf{Tema:} Desglose de cuota en sistema francés. \textbf{Explicación:} El interés del primer periodo se calcula sobre el saldo inicial: $I_1 = 100000 \times 0.10 = 10000.00$ soles. La amortización es la diferencia: $A_1 = Cuota - I_1 = 31547.08 - 10000.00 = 21547.08$ soles.}

\sol{52}{D}{\textbf{Tema:} Sistema de amortización alemán. \textbf{Explicación:} Amortización fija $= 60000/3 = 20000$ soles. Año 1: interés $6000$, cuota $26000$. Año 2: saldo $40000$, interés $4000$, cuota $24000$. Año 3: saldo $20000$, interés $2000$, cuota $22000$.}

\sol{53}{B}{\textbf{Tema:} Precio de bono cupón cero. \textbf{Explicación:} $P = \frac{Valor\ Nominal}{(1 + YTM)^n} = \frac{1000}{(1.08)^2} \approx 857.34$ soles.}

\sol{54}{C}{\textbf{Tema:} Sensibilidad de precios por Duración Modificada. \textbf{Explicación:} $\frac{\Delta P}{P} \approx -D_{mod} \times \Delta y = -4.5 \times (+0.005) = -0.0225 = -2.25\%$. El precio del bono cae un $2.25\%$.}

\sol{55}{B}{\textbf{Tema:} Cálculo de la TIRM. \textbf{Explicación:} $1000(1 + TIRM)^2 = 1460 \implies (1 + TIRM)^2 = 1.46 \implies 1 + TIRM = 1.2083 \implies TIRM = 20.83\%$.}

\sol{56}{B}{\textbf{Tema:} Punto de equilibrio operativo. \textbf{Explicación:} $Q = \frac{Costos\ Fijos}{P - CV} = \frac{120000}{5.0 - 2.0} = \frac{120000}{3.0} = 40000$ metros cúbicos.}

\sol{57}{D}{\textbf{Tema:} Grado de apalancamiento operativo (GAO). \textbf{Explicación:} $\% \Delta EBIT = GAO \times \% \Delta Ventas = 2.5 \times 8.0\% = 20.0\%$.}

\sol{58}{D}{\textbf{Tema:} Grado de apalancamiento total (GAT). \textbf{Explicación:} $GAT = GAO \times GAF = 2.0 \times 1.5 = 3.0$. Mide el impacto porcentual total sobre la utilidad neta por acción ante una variación del 1\% en las ventas.}

\sol{59}{C}{\textbf{Tema:} Método del Valor Presente Ajustado (APV). \textbf{Explicación:} $APV = VAN_{unlevered} + VP(Financiamiento)$. Desagrega el valor intrínseco del proyecto de los efectos colaterales del financiamiento, siendo ideal cuando el nivel de deuda no es una proporción constante del valor de mercado.}

\sol{60}{D}{\textbf{Tema:} Opciones reales de expansión en infraestructura. \textbf{Explicación:} La flexibilidad gerencial para ampliar capacidad si las condiciones de demanda son favorables tiene valor económico real (Opción Call). El enfoque de opciones reales reconoce que la gerencia no está atada pasivamente a un plan rígido inicial.}

\sol{61}{D}{\textbf{Tema:} Modigliani-Miller Proposición II con impuestos. \textbf{Explicación:} $K_e = K_u + (K_u - K_d)(1 - T)(D/E) = 0.10 + (0.10 - 0.06)(0.70)(0.50) = 0.10 + 0.04(0.35) = 0.10 + 0.014 = 0.114 = 11.40\%$.}

\sol{62}{C}{\textbf{Tema:} Consistencia en evaluación con inflación. \textbf{Explicación:} Mezclar flujos nominales con tasas reales (o viceversa) introduce severas distorsiones. La consistencia matemática exige que el numerador y el denominador incorporen (o excluyan) la prima de inflación en idéntica medida.}

\sol{63}{D}{\textbf{Tema:} Valor terminal por perpetuidad. \textbf{Explicación:} $Valor\ Terminal_{10} = \frac{FCLP_{10}}{r} = \frac{80000}{0.10} = 800000$ soles evaluados en el año 10. Para sumarlo al VAN se debe descontar a valor presente: $\frac{800000}{(1.10)^{10}}$.}

\sol{64}{A}{\textbf{Tema:} Cálculo de ratio P/E. \textbf{Explicación:} $P/E = \frac{Precio}{Utilidad\ por\ Acci\acute{o}n} = \frac{30.0}{2.50} = 12.0$ veces. Los inversionistas pagan 12 soles por cada sol de beneficio anual generado por la compañía.}

\sol{65}{A}{\textbf{Tema:} Cálculo de YTM de bono cupón cero. \textbf{Explicación:} $(1 + YTM)^3 = \frac{1000}{751.3146} = 1.331 \implies 1 + YTM = \sqrt[3]{1.331} = 1.10 \implies YTM = 10.00\%$ anual.}

\subsection*{Bloque C. Casos con implicancias regulatorias y sector saneamiento}

\sol{66}{C}{\textbf{Tema:} Estructura de capital regulatoria de SUNASS. \textbf{Explicación:} SUNASS establece una estructura de financiamiento eficiente y normativa (empresa modelo eficiente) para que la tarifa remunere un costo de capital prudente, sin convalidar sobrecostos si la EPS está mal financiada ni desincentivar la inversión si carece de deuda.}

\sol{67}{D}{\textbf{Tema:} Tasa libre de riesgo en regulación. \textbf{Explicación:} Los US Treasuries a largo plazo son el estándar financiero internacional para $R_f$. Dado que la infraestructura de agua tiene vidas útiles de 20 a 40 años, se eligen bonos soberanos con vencimientos largos compatibles con la vida de los activos.}

\sol{68}{B}{\textbf{Tema:} Riesgo país EMBI+ en WACC regulatorio. \textbf{Explicación:} El EMBI+ cuantifica el riesgo soberano de impago y estabilidad macroeconómica de invertir en el Perú frente a EE.UU., sumándose a la tasa libre de riesgo en el modelo CAPM extendido para mercados emergentes.}

\sol{69}{C}{\textbf{Tema:} Definición de Base de Activos Regulatorios. \textbf{Explicación:} La BAR representa el capital físico productivo reconocido por el regulador (plantas, redes, reservorios). La retribución al capital en la tarifa se calcula como $Retribuci\acute{o}n = BAR \times WACC$, asegurando el retorno del capital invertido.}

\sol{70}{A}{\textbf{Tema:} Metodología del Valor Nuevo de Reemplazo (VNR). \textbf{Explicación:} El VNR simula el costo de reproducir el servicio hoy con la mejor tecnología disponible y eficiente. Evita subvaluar activos amortizados contablemente que aún operan y evita sobrevaluar infraestructura obsoleta con tecnología costosa superada.}

\sol{71}{C}{\textbf{Tema:} Justificación del WACC en empresas públicas. \textbf{Explicación:} Si la tarifa solo cubriera costos operativos (OPEX) y depreciación contable con retorno cero, la empresa descapitalizaría sus redes y sería incapaz de financiar nuevas plantas de tratamiento o acceder a préstamos del BID o KfW sin depender de subsidios del MEF.}

\sol{72}{A}{\textbf{Tema:} Evaluación social en Invierte.pe. \textbf{Explicación:} Invierte.pe evalúa si el proyecto genera bienestar social neto para el país. Corrige distorsiones de mercado (impuestos indirectos, subsidios, desempleo) mediante factores de conversión social aplicados a la mano de obra, divisa y combustibles.}

\sol{73}{C}{\textbf{Tema:} Tasa Social de Descuento del MEF. \textbf{Explicación:} La TSD oficial de Invierte.pe se ubica en $8.0\%$ anual real (reducida desde el histórico 10\% y 9\%). Representa el costo de oportunidad social de los fondos públicos utilizados en proyectos de inversión pública en el Perú.}

\sol{74}{D}{\textbf{Tema:} Factor de corrección de mano de obra no calificada. \textbf{Explicación:} En regiones con subempleo, emplear a un trabajador rural en una obra de saneamiento no reduce la producción agrícola en la misma cuantía del salario pagado. El factor de corrección descuenta esa brecha para reflejar su menor costo social real.}

\sol{75}{B}{\textbf{Tema:} Análisis Costo-Eficacia en Invierte.pe. \textbf{Explicación:} El Análisis Costo-Eficacia se utiliza cuando los beneficios sociales son evidentes pero no monetizables con rigor econométrico. Se fija el objetivo social (ej. 100\% de cobertura y agua clorada) y se elige la alternativa con menor Ratio Costo-Eficacia ($RCE$).}

\sol{76}{D}{\textbf{Tema:} Ratio DSCR en contratos de crédito de EPS. \textbf{Explicación:} El DSCR (Debt Service Coverage Ratio) mide la holgura de caja de la empresa prestadora para pagar amortización e intereses con el flujo generado por las tarifas. Si cae por debajo de 1.20, los bancos restringen nuevos desembolsos.}

\sol{77}{D}{\textbf{Tema:} Cuentas por cobrar del sector público en EPS. \textbf{Explicación:} La morosidad de entidades del Estado con las EPS genera un grave problema de capital de trabajo: las EPS deben pagar energía eléctrica, cloro y planillas al contado, mientras sus cuentas por cobrar públicas tardan meses en liquidarse.}

\sol{78}{C}{\textbf{Tema:} Depreciación regulatoria y fondo de reposición. \textbf{Explicación:} A diferencia de una empresa comercial donde la depreciación es un mero asiento contable de costo hundido, en la regulación de SUNASS la depreciación alimenta cuentas bancarias restringidas de inversión para garantizar la reposición continua de la infraestructura.}

\sol{79}{D}{\textbf{Tema:} Estructura RPI y RPMO en concesiones de saneamiento. \textbf{Explicación:} El esquema RPI/RPMO separa el repago del capital invertido (compromiso firme de pago del Estado o de la tarifa) del costo operativo mensual, el cual se penaliza si la planta no cumple con los límites máximos permisibles de vertimiento.}

\sol{80}{A}{\textbf{Tema:} Project Finance en infraestructura hídrica. \textbf{Explicación:} En Project Finance los prestamistas no tienen recurso contra los accionistas promotores si el proyecto quiebra: el riesgo descansa en la viabilidad intrínseca del flujo de caja del contrato de concesión y en la solidez del marco regulatorio tarifario.}

\sol{81}{A}{\textbf{Tema:} Pérdida del escudo fiscal en empresas insolventes. \textbf{Explicación:} El escudo fiscal de los intereses solo genera valor en efectivo si la empresa tiene utilidades tributarias gravables contra las cuales deducir el gasto. Sin ganancias operativas netas, el beneficio del escudo impositivo de Modigliani-Miller desaparece.}

\sol{82}{D}{\textbf{Tema:} Intangibilidad de fondos MERESE. \textbf{Explicación:} La regulación de SUNASS exige estricta intangibilidad contable y financiera: los fondos de la cuota MERESE y GRD no pueden mezclarse con el flujo operativo corriente de la EPS ni destinarse a gastos administrativos generales.}

\sol{83}{A}{\textbf{Tema:} Descalce cambiario en empresas de servicios públicos. \textbf{Explicación:} Como las tarifas de agua no están dolarizadas, un salto en el tipo de cambio encarece el pago de amortización e intereses en moneda nacional, comprometiendo la liquidez y solvencia de la prestadora si no cuenta con cobertura cambiaria (hedging).}

\sol{84}{A}{\textbf{Tema:} Opciones reales modulares en saneamiento. \textbf{Explicación:} La construcción modular es una opción real de crecimiento: evita hundir capital en megaobras con capacidad ociosa durante décadas y permite adaptar la inversión según la evolución real de la demanda y el cambio climático.}

\sol{85}{D}{\textbf{Tema:} Causales financieras de ingreso al RAT de OTASS. \textbf{Explicación:} El Decreto Legislativo 1280 establece causales objetivas de ingreso al RAT: insolvencia técnica y financiera (patrimonio neto menor a cero o pérdidas de más de dos tercios del capital), quiebra operativa y grave incapacidad de gestión del prestador municipal.}

\sol{86}{B}{\textbf{Tema:} Obras por Impuestos (OxI) en saneamiento. \textbf{Explicación:} OxI agiliza la inversión pública: una empresa privada adelanta el capital y ejecuta la obra de agua y alcantarillado, recibiendo del MEF Certificados de Inversión Pública Regional y Local (CIPRL) aplicables al pago de su impuesto a la renta.}

\sol{87}{A}{\textbf{Tema:} Contingencias ambientales y pasivos ocultos. \textbf{Explicación:} No provisionar contingencias por descargas de aguas residuales sin tratamiento (sancionables por OEFA y ANA) infla artificialmente el patrimonio de la EPS, engañando a auditores y retrasando la inversión necesaria en remediación ambiental.}

\sol{88}{A}{\textbf{Tema:} Evaluación por CAE en micromedición. \textbf{Explicación:} La comparación entre tecnologías de medición con distinta durabilidad y costos exige el Costo Anual Equivalente neto del beneficio de reducción de ANF aparente. La tecnología con menor costo neto anualizado maximiza la rentabilidad operativa de la EPS.}

\sol{89}{C}{\textbf{Tema:} Fundamento financiero del periodo regulatorio quinquenal. \textbf{Explicación:} Las inversiones en infraestructura hídrica tardan varios años en ejecutarse y madurar. Un quinquenio brinda certidumbre regulatoria a la empresa para invertir y capturar ganancias de eficiencia antes de que el regulador reajuste la tarifa base en el siguiente ciclo.}

\sol{90}{D}{\textbf{Tema:} Equilibrio económico-financiero en concesiones. \textbf{Explicación:} Las cláusulas de equilibrio económico-financiero blindan la bancabilidad de los contratos a largo plazo: si el Estado altera ex-post las condiciones regulatorias o tributarias de forma sustancial (hecho del príncipe), el contrato prevé ajustes en la tarifa o plazos para restituir el VAN original pactado.}

\subsection*{Bloque D. Ampliación: Teoría y Aplicaciones Prácticas (91 a 180)}

\sol{91}{B}{Por el límite notable $\lim_{m\to\infty}(1 + r/m)^m = e^r$, la tasa efectiva anual bajo capitalización continua es exactamente $TEA = e^r - 1$.}
\sol{92}{C}{Por Fisher, $1 + r_{ex-post} = \frac{1 + i}{1 + \pi_{efectiva}}$. Si $\pi_{efectiva} > \pi^e$, la tasa real efectiva cae y el deudor paga con moneda de menor poder adquisitivo.}
\sol{93}{D}{El crecimiento solo crea valor ($PVGO > 0$) cuando los proyectos reinvertidos generan $VAN > 0$, es decir, $ROE > k_e$. Si $ROE = k_e$, el $PVGO = 0$.}
\sol{94}{A}{La fórmula estándar de anualidad vencida coloca el valor presente un periodo antes del primer pago (en $t=k$); luego se trae a $t=0$ dividiendo entre $(1+r)^k$.}
\sol{95}{C}{La teoría de preferencia por la liquidez añade una prima positiva creciente con el plazo para compensar el mayor riesgo de precio de bonos largos, inclinando la curva hacia arriba.}
\sol{96}{D}{Sin arbitraje, invertir a 2 años a la tasa $s_2$ debe rendir igual que invertir a 1 año a $s_1$ y pactar la reinversión del año 2 a $f_{1,2}$: $(1+s_2)^2 = (1+s_1)(1+f_{1,2})$.}
\sol{97}{A}{En el sistema francés la cuota $C$ es constante; al reducirse el saldo deudor, los intereses bajan y la amortización sube. En el alemán la amortización es fija, bajando el saldo más rápido (menos intereses totales).}
\sol{98}{B}{A mayor $m$, los intereses se capitalizan antes y generan intereses sobre intereses, elevando la TEA de forma monótona y cóncava hacia su cota superior $e^{TNA}-1$.}
\sol{99}{D}{La utilidad neta se rige por el devengado (registra facturación aun si no se cobra) y no refleja desembolsos de capital (CAPEX) ni amortizaciones de principal de deuda.}
\sol{100}{A}{El FCL pertenece a todos los proveedores de capital (acreedores y accionistas), por lo que se descuenta al WACC; el FCA es residual para accionistas tras deuda y se descuenta a $k_e$.}
\sol{101}{B}{Aumentar el CTN (más inventarios o cuentas por cobrar netas) inmoviliza caja ($-\Delta CTN$ en el FCL); al finalizar el proyecto ese capital de trabajo se liquida y se recupera como flujo positivo.}
\sol{102}{C}{Al deducirse antes de impuestos y luego sumarse de vuelta al NOPAT, la depreciación incrementa el flujo de caja libre en exactamente $T \cdot Dep$ (escudo fiscal de la depreciación).}
\sol{103}{A}{Como $EVA = (ROIC - WACC) \times \text{Capital Invertido}$, si $ROIC < WACC$ la empresa no cubre el costo de oportunidad de su capital y destruye valor económico ($EVA < 0$).}
\sol{104}{B}{La identidad DuPont de 5 pasos expresa $ROE = \frac{UN}{EBT} \times \frac{EBT}{EBIT} \times \frac{EBIT}{\text{Ventas}} \times \frac{\text{Ventas}}{\text{Activos}} \times \frac{\text{Activos}}{\text{Patrimonio}}$.}
\sol{105}{C}{Por la relación $ROE = ROIC + (ROIC - k_d(1-T))\frac{D}{E}$, endeudarse eleva el $ROE$ si y solo si los activos rinden más que el costo neto de la deuda ($ROIC > k_d(1-T)$).}
\sol{106}{D}{$GAO = \frac{\%\Delta EBIT}{\%\Delta Q}$ (riesgo por costos fijos operativos) y $GAF = \frac{\%\Delta UPA}{\%\Delta EBIT}$ (riesgo por gastos financieros fijos); además $GAT = GAO \times GAF$.}
\sol{107}{B}{El uso de un activo existente tiene un costo de oportunidad igual a su mejor uso alternativo (valor de mercado neto de impuestos hoy), a diferencia de los costos hundidos irrecuperables.}
\sol{108}{C}{Con dos cambios de signo ($-, +, \dots, +, -$), la ecuación $VAN(r)=0$ puede tener hasta 2 raíces reales positivas (múltiples TIR). El VAN y la TIRM evitan esta ambigüedad.}
\sol{109}{D}{Como $VAN_A(r) - VAN_B(r) = VAN_{A-B}(r) = 0$, la tasa de Fisher $r_F$ es exactamente la TIR del flujo incremental $(FC_A - FC_B)$.}
\sol{110}{A}{Con presupuesto limitado en un periodo y proyectos divisibles, ordenar por $IR$ (o razón $VAN/I_0$) maximiza el VAN total por unidad de capital escaso (si son indivisibles, se usa programación entera).}
\sol{111}{C}{Un activo que dura 4 años deberá reemplazarse antes que uno de 6 años; el CAE anualiza el VAC a lo largo de la vida útil (equivalente a igualar horizontes mediante cadena de reemplazos).}
\sol{112}{D}{Se reemplaza cuando el costo incremental de operar el activo viejo un año más (O\&M + pérdida de valor de rescate + costo de oportunidad del rescate) excede el CAE del activo nuevo.}
\sol{113}{A}{Aunque descuenta los flujos al costo de capital, el Payback descontado descarta todo flujo posterior al periodo de corte, pudiendo rechazar proyectos de gran VAN con flujos tardíos.}
\sol{114}{B}{El análisis de sensibilidad varía una variable a la vez (*ceteris paribus*); Monte Carlo modela distribuciones conjuntas y correlaciones para obtener la distribución de probabilidad completa del VAN.}
\sol{115}{D}{El VAN expandido equivale a $VAN_{estático} + \text{Valor de Opciones Reales}$; el VAN tradicional omite la capacidad de diferir, expandir o abandonar conforme se resuelve la incertidumbre.}
\sol{116}{A}{Como $\sigma_p^2 = w_1^2\sigma_1^2 + w_2^2\sigma_2^2 + 2w_1 w_2 \rho_{12}\sigma_1\sigma_2$, para cualquier $\rho_{12} < 1$ se cumple $\sigma_p < w_1\sigma_1 + w_2\sigma_2$, existiendo ganancia de diversificación.}
\sol{117}{B}{La decisión de inversión (elegir el portafolio tangente de máximo Sharpe) se separa de la decisión de financiamiento/riesgo (combinar ese portafolio con el activo libre de riesgo sobre la CML).}
\sol{118}{C}{La CML relaciona $E[R_p]$ con $\sigma_p$ (pendiente = ratio de Sharpe del mercado) solo para portafolios eficientes; la SML relaciona $E[R_i]$ con $\beta_i$ (pendiente = $E[R_m]-r_f$) para todo activo en equilibrio.}
\sol{119}{A}{$\alpha_i = E[R_i] - [r_f + \beta_i(E[R_m]-r_f)] > 0$ significa que el activo promete más retorno del exigido por su beta: está subvaluado (barato), atrayendo compras que elevan su precio y bajan su retorno hacia la SML.}
\sol{120}{B}{Sharpe divide el exceso de retorno por el riesgo total ($\sigma_p$); Treynor lo divide por el riesgo sistemático ($\beta_p$), apropiado cuando el riesgo idiosincrático ya está diversificado.}
\sol{121}{C}{Hamada separa riesgo operativo ($\beta_U$) del financiero: primero se desapalanca cada comparable ($\beta_U = \frac{\beta_L}{1+(1-T)D/E}$) y luego se reapalanca con la $D/E$ y $T$ regulatorias.}
\sol{122}{D}{Sin impuestos, $k_e = k_u + (k_u - k_d)(D/E)$. Sustituyendo en $WACC = \frac{E}{V}k_e + \frac{D}{V}k_d$, los efectos se cancelan exactamente y $WACC = k_u$ permanece constante.}
\sol{123}{B}{En el modelo Trade-off, $V_L = V_U + VP(\text{Escudo Fiscal}) - VP(\text{Costos de Quiebra})$; el óptimo interior $D^*$ iguala el beneficio marginal tributario al costo marginal de insolvencia.}
\sol{124}{C}{Por selección adversa, emitir acciones señala al mercado que la gerencia las cree sobrevaluadas (haciendo caer el precio); por ello se prefieren fondos internos, luego deuda y al final acciones.}
\sol{125}{D}{Con responsabilidad limitada, el patrimonio actúa como una opción *call* sobre los activos: cerca de la quiebra, mayor volatilidad aumenta el valor del patrimonio a expensas de los acreedores.}
\sol{126}{A}{Sin fricciones, si un accionista desea efectivo cuando la firma no paga dividendo, vende una fracción de sus acciones ('dividendo casero'), dejando intacta su riqueza total.}
\sol{127}{C}{El WACC constante asume que $D/V$ se rebalancea continuamente. El APV ($VAN_{desapalancado} + VP(\text{efectos financieros})$) maneja con exactitud calendarios de deuda fijos, subsidios o *Project Finance*.}
\sol{128}{D}{Cuando la tasa exigida por el mercado supera la tasa cupón ($YTM > c$), el bono debe ofrecer una ganancia de capital para compensar, cotizando por debajo de su valor nominal (bajo par).}
\sol{129}{A}{La duración es el promedio ponderado temporal de los flujos: cupones altos adelantan el peso temporal (baja $D$); sin cupones (cupón cero), todo el flujo ocurre en $n$, por lo que $D = n$.}
\sol{130}{B}{La recta tangente (Duración Modificada) yace siempre por debajo de la curva convexa precio-YTM; por tanto, siempre subestima el precio real (subestima la ganancia si baja $y$ y sobreestima la pérdida si sube $y$).}
\sol{131}{D}{Cuando la duración de la cartera iguala el horizonte $H$, el efecto precio y el efecto reinversión ante un cambio pequeño y paralelo en las tasas se compensan mutuamente en $t = H$.}
\sol{132}{A}{En $t=T$, ambos portafolios pagan exactamente $\max(S_T, K)$. Por ley de un solo precio (no arbitraje), deben costar exactamente lo mismo hoy en $t=0$.}
\sol{133}{B}{Las opciones tienen pagos asimétricos (pérdida acotada a la prima y ganancia potencial ilimitada); mayor volatilidad aumenta la probabilidad de pagos favorables, elevando tanto *call* como *put* ($\text{Vega} > 0$).}
\sol{134}{C}{Un *Forward* es una obligación simétrica sin prima inicial; una *Call* otorga el derecho (no obligación) de comprar a $K$, protegiendo contra alzas del dólar y permitiendo beneficiarse si baja, a cambio de una prima.}
\sol{135}{A}{Esperar para invertir equivale a una *call* (derecho a pagar $I$ para adquirir el proyecto); abandonar por el valor de salvamento $S$ equivale a una *put* (derecho a vender el proyecto al precio $S$).}
\sol{136}{B}{SUNASS fija la tarifa media igualando el valor presente de ingresos proyectados al valor presente de costos eficientes (O\&M, inversiones financiadas por tarifa e impuestos) descontados al WACC real.}
\sol{137}{C}{Al no cotizar en bolsa y tener estructuras financieras distorsionadas, SUNASS emplea betas de empresas comparables (sector agua/utilities) ajustados por Hamada a una estructura de capital objetivo eficiente.}
\sol{138}{D}{El VNR valora la infraestructura eficiente requerida hoy a precios actuales de reposición, dando señales económicas correctas frente a la distorsión inflacionaria de los costos históricos.}
\sol{139}{B}{Los activos donados no implican costo de oportunidad de capital para la EPS (no se remuneran con WACC para no cobrar dos veces al ciudadano), pero sí requieren cubrir O\&M y futura reposición.}
\sol{140}{C}{El PMO planifica infraestructura y oferta-demanda en un horizonte de largo plazo (30 años), sirviendo de base para determinar la fórmula tarifaria, estructura tarifaria y metas de gestión cada 5 años.}
\sol{141}{D}{Si la tasa permitida supera el costo de capital, la regulación por tasa de retorno incentiva sobreinversión en capital (Averch-Johnson) y traslado de costos; *Price Cap* incentiva eficiencia en costos.}
\sol{142}{A}{La evaluación social mide el impacto neto en el bienestar del país: usa precios sombra (sociales), incorpora externalidades (ej. salud pública) y descuenta los flujos a la TSD fijada por el MEF.}
\sol{143}{C}{Desde la perspectiva del país en su conjunto, un impuesto o subsidio es una transferencia entre privados y el Estado; el costo económico real corresponde a los recursos netos de distorsiones fiscales.}
\sol{144}{D}{Donde hay subempleo, el costo de oportunidad social de la mano de obra no calificada es inferior al salario de mercado ($FC < 1$). Los impuestos/aranceles hacen que el valor social de una divisa supere al tipo de cambio de mercado ($FC > 1$).}
\sol{145}{A}{El beneficio social de pasar de cisterna (alto costo) a red pública (menor costo) es el área bajo la curva de demanda: recursos liberados del consumo antiguo más excedente por consumo incremental.}
\sol{146}{B}{Invierte.pe usa Costo-Beneficio cuando los beneficios son monetizables (agua potable); usa Costo-Eficacia (menor costo social por beneficiario o $m^3$ tratado) cuando monetizar beneficios resulta complejo (alcantarillado/PTAR).}
\sol{147}{D}{Una TSD de $8\%$ castiga drásticamente beneficios intergeneracionales (a 50 años, $1.08^{-50} \approx 0.02$). Por equidad intergeneracional e incertidumbre de largo plazo, el MEF establece tasas sociales menores/decrecientes para proyectos ambientales.}
\sol{148}{A}{El *Value for Money* compara el costo total del ciclo de vida (inversión, O\&M y riesgos transferidos) entre APP y obra pública tradicional (Comparador Público-Privado), buscando eficiencia neta.}
\sol{149}{B}{Transferir riesgos incontrolables al privado encarece innecesariamente el contrato (prima de riesgo excesiva). La regla óptima asigna cada riesgo a quien mejor puede controlarlo al menor costo.}
\sol{150}{C}{En esquemas de pago por disponibilidad, el concesionario recibe el RPI/RPMO (respaldado usualmente por fideicomiso) sujeto a cumplir los indicadores de disponibilidad y calidad (descuentos por incumplimiento).}
\sol{151}{A}{En OxI, el privado adelanta el financiamiento y ejecución de un proyecto público viable (Invierte.pe) y recibe CIPRL/CIPGN que aplica contra su Impuesto a la Renta, descontándose luego del canon/recursos de la entidad.}
\sol{152}{B}{Un *Working Ratio* $< 1$ (idealmente $< 0.70$) implica que los ingresos operativos cubren con holgura los costos operativos en efectivo, dejando excedente (margen EBITDA positivo) para servicio de deuda e inversiones.}
\sol{153}{C}{La intangibilidad obliga a depositar los porcentajes tarifarios asignados a MERESE/GRD en cuentas exclusivas; su desvío a gasto corriente está prohibido y sujeto a sanción por SUNASS.}
\sol{154}{D}{$TEA = \left(1 + \frac{0.12}{4}\right)^4 - 1 = (1.03)^4 - 1 = 1.125509 - 1 = 12.55\%$.}
\sol{155}{B}{Bajo capitalización continua: $VF = VP \cdot e^{rt} = 1{,}000 \cdot e^{0.08 \times 5} = 1{,}000 \cdot e^{0.40} \approx S/\ 1{,}491.82$.}
\sol{156}{C}{Por Fisher exacto: $1 + r = \frac{1 + i}{1 + \pi} = \frac{1.1025}{1.05} = 1.05 \Rightarrow r = 5.00\%$.}
\sol{157}{D}{$VP = 1{,}000 \times \frac{1 - (1.10)^{-3}}{0.10} = 1{,}000 \times \frac{1 - 0.751315}{0.10} = S/\ 2{,}486.85$.}
\sol{158}{A}{Primero hallamos $D_1 = D_0(1+g) = 4.00(1.05) = S/\ 4.20$. Luego, $P_0 = \frac{D_1}{k_e - g} = \frac{4.20}{0.13 - 0.05} = \frac{4.20}{0.08} = S/\ 52.50$.}
\sol{159}{C}{$g = ROE \times b = 0.15 \times 0.40 = 6.0\%$. El dividendo esperado es $D_1 = 0.60 \times 10 = S/\ 6.00$, y $P_0 = \frac{6.00}{0.12 - 0.06} = S/\ 100.00$.}
\sol{160}{D}{$CTN = AC - PC = 450 - 300 = S/\ 150$ millones; Razón Corriente $= \frac{AC}{PC} = \frac{450}{300} = 1.50$.}
\sol{161}{A}{La Prueba Ácida excluye los inventarios: $\frac{AC - \text{Inventarios}}{PC} = \frac{320 - 80}{200} = \frac{240}{200} = 1.20$.}
\sol{162}{B}{$ROE = \text{Margen Neto} \times \text{Rotación} \times \text{Multiplicador} = 0.08 \times 1.25 \times 1.80 = 0.18 = 18.0\%$.}
\sol{163}{D}{$FCL = EBIT(1-T) + Dep - CAPEX - \Delta CTN = 200(0.70) + 50 - 80 - 20 = 140 + 50 - 100 = S/\ 90$ mil.}
\sol{164}{A}{$EVA = NOPAT - (WACC \times \text{Capital Invertido}) = 65 - (0.10 \times 500) = 65 - 50 = S/\ 15$ millones.}
\sol{165}{B}{$VAN = -1{,}000 + \frac{550}{1.10} + \frac{726}{(1.10)^2} = -1{,}000 + 500 + 600 = S/\ 100.00$.}
\sol{166}{C}{El Índice de Rentabilidad es $IR = \frac{VP(\text{Flujos})}{I_0} = \frac{1{,}320}{1{,}000} = 1.32$ (como $IR > 1$, el proyecto crea valor).}
\sol{167}{A}{Resolviendo $-100 + \frac{60}{1+r} + \frac{72}{(1+r)^2} = 0$ con $x = 1+r$: $25x^2 - 15x - 18 = (5x-6)(5x+3) = 0 \Rightarrow x = 1.20 \Rightarrow TIR = 20.0\%$.}
\sol{168}{B}{Igualando $VAN_A = VAN_B$ (flujo incremental $B-A$: $0, -110, +121$): $-\frac{110}{1+r_F} + \frac{121}{(1+r_F)^2} = 0 \Rightarrow 1 + r_F = \frac{121}{110} = 1.10 \Rightarrow r_F = 10.0\%$.}
\sol{169}{C}{$CAE = VAC \times \frac{r}{1 - (1+r)^{-n}} = 1{,}000 \times \frac{0.10}{1 - 0.7513} = \frac{100}{0.2487} \approx S/\ 402.11$.}
\sol{170}{D}{Los valores presentes son $VP_1 = 500$, $VP_2 = 400$ (acumulado $900$) y $VP_3 = 200$. Faltan $S/\ 100$ en el año 3: $2 + \frac{100}{200} = 2.50$ años.}
\sol{171}{B}{Aplicando el CAPM emergente: $k_e = r_f + \beta_L(E[R_m] - r_f) + RP = 4.0\% + 0.80(6.0\%) + 2.0\% = 10.80\%$.}
\sol{172}{C}{$\beta_L = \beta_U\left[1 + (1 - T)\frac{D}{E}\right] = 0.60[1 + (1 - 0.30)(0.50)] = 0.60(1.35) = 0.81$.}
\sol{173}{D}{$WACC = \frac{E}{V}k_e + \frac{D}{V}k_d(1-T) = 0.60(12.0\%) + 0.40(8.0\%)(0.75) = 7.20\% + 2.40\% = 9.60\%$.}
\sol{174}{A}{Según MM I con impuestos corporativos y deuda perpetua: $V_L = V_U + T \cdot D = 800 + 0.30(200) = S/\ 860$ millones.}
\sol{175}{C}{Por MM II sin impuestos: $k_e = k_u + (k_u - k_d)\frac{D}{E} = 11.0\% + (11.0\% - 6.0\%)(0.80) = 11.0\% + 4.0\% = 15.0\%$.}
\sol{176}{D}{$P_B = \frac{80}{1.10} + \frac{1{,}080}{(1.10)^2} = 72.73 + 892.56 = S/\ 965.29$ (cotiza bajo par porque $c < YTM$).}
\sol{177}{A}{Primero, $D_{mod} = \frac{5.40}{1.08} = 5.00$. Luego, $\frac{\Delta P}{P} \approx -D_{mod}\Delta y = -5.00(0.005) = -0.025 = -2.50\%$.}
\sol{178}{B}{Por paridad: $P = C + \frac{K}{1+r} - S_0 = 12.00 + \frac{110.25}{1.05} - 100.00 = 12.00 + 105.00 - 100.00 = S/\ 17.00$.}
\sol{179}{D}{Retorno CAPM esperado: $4.0\% + 1.25(10.0\% - 4.0\%) = 11.5\%$, luego $\alpha_p = 14.0\% - 11.5\% = +2.5\%$. El Ratio de Sharpe es $SR = \frac{14.0\% - 4.0\%}{20.0\%} = 0.50$.}
\sol{180}{A}{Aplicando cada factor de corrección: $600(0.85) + 200(0.50) + 200(1.10) = 510 + 100 + 220 = S/\ 830$ mil.}
